\begin{table}[htbp]
\centering
\begingroup
\small
\setlength{\tabcolsep}{4pt}
\begin{tabular}{llll}
\toprule
Split & Healthy & Outer race & Inner race \\
\midrule
fold A & K001, K002, K003 & KA04, KA15, KA16 & KI04, KI14, KI16 \\
fold B & K004, K005, K006 & KA22, KA30 & KI17, KI18, KI21 \\
\midrule
S00 & K001, K003, K004 & KA16, KA22, KA30 & KI04, KI16, KI18 \\
S01 & K002, K004, K006 & KA15, KA22, KA30 & KI04, KI17, KI21 \\
S02 & K002, K003, K005 & KA15, KA16, KA30 & KI16, KI17, KI21 \\
S03 & K002, K003, K004 & KA04, KA16, KA22 & KI16, KI18, KI21 \\
S04 & K001, K002, K006 & KA15, KA16, KA30 & KI04, KI14, KI18 \\
S05 & K003, K005, K006 & KA04, KA16, KA30 & KI04, KI17, KI21 \\
S06 & K002, K003, K006 & KA16, KA22, KA30 & KI04, KI17, KI18 \\
S07 & K003, K005, K006 & KA04, KA16, KA22 & KI04, KI14, KI16 \\
S08 & K001, K002, K006 & KA15, KA22, KA30 & KI04, KI16, KI21 \\
S09 & K002, K003, K005 & KA04, KA15, KA16 & KI16, KI17, KI21 \\
\bottomrule
\end{tabular}\par
\endgroup
\caption{Paderborn bearing splits. Each row lists the source bearings of one Paderborn split. The clean target of a split is the rest of the real-damage pool of 6 healthy, 5 outer-race and 6 inner-race bearings, and the leaky target is the source bearings at the target condition. Fold~B has two outer-race source bearings because the pool holds five.}
\label{tab:splits}
\end{table}
